\section{Closing the simultaneous induction}
\begin{theorem}\label{sharpclosure}
For every faithful transitive string C-group of degree $n$, other than
$A_n,S_n$, the rank is at most $n/2+1$. In the imprimitive case with
$n\geq10$, equality requires two-point blocks with an odd number of
blocks, and has exactly the markings classified above. In the primitive
case the rank is strictly below $n/2$ except for the explicitly enumerated
degree-six and degree-ten markings.
\end{theorem}
\begin{proof}
We prove these assertions simultaneously by strong induction on $n$.
The small imprimitive actions of degrees six, eight and nine are covered
by the complete required census. It searches through rank $n-2$, the
Whiston bound for a proper subgroup other than $S_n$, without using any
imprimitive rank theorem. A prime degree has no nontrivial block system.
For degree four the possible imprimitive groups are subgroups of $D_8$;
its Frattini quotient is $C_2^2$, so any independent generating family
has at most two members. The primitive checks through degree twenty-five
are independent of the desired induction assertion. These observations
supply every small endpoint needed below.

At degree $n$, the maximal-block argument handles $k,m>2$. For blocks
of size two it uses only the previously accepted independent-length
bounds on primitive quotient images, the explicit binary-module lemma,
and the complete extremal $S_5,S_6$ checks. It does not apply a C-group
rank theorem to a quotient that might have lost IP. For two blocks, the
finite ten-action lemma treats $5\leq k\leq8$. When $k\geq9$, the
normalized $H$ is itself a string C-group of degree $k<n$, so the
induction hypothesis supplies precisely the assumption in
Theorem~\ref{twoblocks}. Thus the whole imprimitive conclusion at degree
$n$ follows with its complete equality classification.

For a primitive group, the product-action alternative of Mar\'oti has
a faithful imprimitive action of degree $m\ell<n$. The induction
hypothesis, not a same-degree result, applies there. The other primitive
alternatives use Whiston, the retained Conder theorem, or the independent
bounded catalogue checks, as proved above. This proves the primitive
conclusion at degree $n$. Every invocation of the asserted sharp rank
bound has therefore decreased the faithful permutation degree, closing
the induction without using the older general imprimitive
independent-length theorem.
\end{proof}
